% Generated by Listas/make_statements.py from pratica-lecture-02.md, item 1. Do not edit.
\textbf{Source.} Muñoz-Garcia \& Toro-Gonzalez (2019), \emph{Strategy and Game Theory: Practice Exercises with Answers}, Ch. 2, Exercise 2 ``Battle of the Sexes'', p.~41 {[}PDF p.~54{]}. Answer in the source: yes, a worked answer follows the exercise. Reproduced verbatim (statement only; the book's worked answer follows it).



\begin{quote}

A husband and a wife are leaving work, and do not remember which event they are attending to tonight. Both of them, however, remember that last night's argument was about either attending to the football game (the most preferred event for the husband) or the opera (the most preferred event for the wife). To make matters worse, their cell phones broke, so they cannot call each other to confirm which event they are attending to. As a consequence, they must simultaneously and independently decide whether to attend to the football game or the opera.



The payoff matrix in Fig. 2.4 describes the preference of the husband (wife) for the football game (opera, respectively), but also indicates that both players prefer to be together rather than being alone (even if they are alone at their most preferred event). Find the set of Nash Equilibria of this game.



{\def\LTcaptype{none} % do not increment counter

\begin{longtable}[]{@{}lll@{}}

\toprule\noalign{}

\emph{Housband} {[}sic{]} ~\emph{Wife} & \emph{Football} & \emph{Opera} \\

\midrule\noalign{}

\endhead

\bottomrule\noalign{}

\endlastfoot

\emph{Football} & 3, 1 & 0, 0 \\

\emph{Opera} & 0, 0 & 1, 3 \\

\end{longtable}

}



\emph{Fig. 2.4} Battle of the sexes game (Normal-form representation)

\end{quote}

